% GENERATED FILE. Edit canonical lessons, then run npm run generate:books.
% canonical-source-hash: fnv1a64:3b9e71417959c77d
% canonical-lessons: TA-C115-nulakam, TA-C115-vanki, TA-C115-cantai, TA-C115-katarkarai, TA-C115-aluvalakam

\chapter{Library, Bank, Market, Beach, Office}
\label{ch:ta-library-bank-market-beach-office}
\hlchaptermodality{115}

\begin{chapteropening}
\textbf{By the end of this chapter:} \emph{I can say a library, a bank, a market, a beach and an office.}

Complete the last lesson of chapter 115: i can say a library, a bank, a market, a beach and an office.
\end{chapteropening}

\section[\texorpdfstring{\ta{நூலகம்} (nūlakam)}{நூலகம் (nūlakam)}]{\ta{நூலகம்} (nūlakam) — a library}

\label{lesson:TA-C115-nulakam}



\begin{quote}
Before the new one: say the Tamil for a park, then the Tamil for a temple.
\end{quote}

\subsection*{You'll want to know: \ta{நூலகம்}}
\textbf{\ta{நூலகம்}} — \emph{nūlakam} — \textquotedblleft{}a library\textquotedblright{}.

\textbf{In the family, and next door}

\noindent
\begin{tabularx}{\linewidth}{@{}>{\raggedright\arraybackslash}X>{\raggedright\arraybackslash}X@{}}
\toprule
\textbf{} & \textbf{word} \\
\midrule
Kannada & \ka{ಗ್ರಂಥಾಲಯ} (\emph{granthālaya}) \\
Telugu & \te{గ్రంథాలయం} (\emph{grandhālayaṁ}) \\
Malayalam & \ml{ഗ്രന്ഥശാല} (\emph{granthaśāla}) \\
Hindi (neighbour) & \dv{पुस्तकालय} (\emph{pustakālay}) \\
English & a library \\
\bottomrule
\end{tabularx}

\begin{grammarlens}[title={What you've built}]
One more word for this chapter.
\end{grammarlens}

\subsection*{Your turn}
\begin{itemize}
\raggedright
  \item \emph{Say it:} \emph{nūlakam}
  \item \emph{Say it:} \emph{nūlakam}, once more
  \item \emph{Say it:} say \emph{nūlakam}
  \item \emph{Recall:} say the Tamil for a park, then the Tamil for a temple, then say \emph{nūlakam} again
\end{itemize}

\begin{tcolorbox}[breakable,colback=teal!4,colframe=teal!35!black,title={Before you move on}]
What does \ta{நூலகம்} mean? (\textquotedblleft{}A library\textquotedblright{}.) Say it once more.
\end{tcolorbox}
\section[\texorpdfstring{\ta{வங்கி} (vaṅki)}{வங்கி (vaṅki)}]{\ta{வங்கி} (vaṅki) — a bank}

\label{lesson:TA-C115-vanki}



\begin{quote}
Before the new one: say the Tamil for a temple, then the Tamil for a library.

Say \textbf{\ta{திங்கள்}} through \textbf{\ta{ஞாயிறு}}. Each day ends in \emph{kiḻamai}; some days carry Tamil words, some Sanskrit ones.
\end{quote}

\subsection*{You'll want to know: \ta{வங்கி}}
\textbf{\ta{வங்கி}} — \emph{vaṅki} — \textquotedblleft{}a bank\textquotedblright{}.

\textbf{In the family, and next door}

\noindent
\begin{tabularx}{\linewidth}{@{}>{\raggedright\arraybackslash}X>{\raggedright\arraybackslash}X>{\raggedright\arraybackslash}X@{}}
\toprule
\textbf{} & \textbf{word} & \textbf{same root} \\
\midrule
Kannada & \ka{ಬ್ಯಾಂಕು} (\emph{byāṅku}) & yes \\
Telugu & \te{బ్యాంకు} (\emph{byāṅku}) & yes \\
Malayalam & \ml{ബാങ്ക്} (\emph{bāṅkŭ}) & yes \\
Hindi (neighbour) & \dv{बैंक} (\emph{baiṅk}) & yes \\
English & a bank &  \\
\bottomrule
\end{tabularx}

\begin{grammarlens}[title={What you've built}]
One more word for this chapter.
\end{grammarlens}

\subsection*{Your turn}
\begin{itemize}
\raggedright
  \item \emph{Say it:} \emph{vaṅki}
  \item \emph{Say it:} \emph{vaṅki}, once more
  \item \emph{Say it:} say \emph{vaṅki}
  \item \emph{Recall:} say the Tamil for a temple, then the Tamil for a library, then say \emph{vaṅki} again
\end{itemize}

\begin{tcolorbox}[breakable,colback=teal!4,colframe=teal!35!black,title={Before you move on}]
What does \ta{வங்கி} mean? (\textquotedblleft{}A bank\textquotedblright{}.) Say it once more.
\end{tcolorbox}
\section[\texorpdfstring{\ta{சந்தை} (cantai)}{சந்தை (cantai)}]{\ta{சந்தை} (cantai) — a market}

\label{lesson:TA-C115-cantai}



\begin{quote}
Before the new one: say the Tamil for a library, then the Tamil for a bank.

One slot changes: \textbf{\ta{இங்கே}}, \textbf{\ta{அங்கே}}, \textbf{\ta{எங்கே}}. Here, there, where?
\end{quote}

\subsection*{You'll want to know: \ta{சந்தை}}
\textbf{\ta{சந்தை}} — \emph{cantai} — \textquotedblleft{}a market\textquotedblright{}.

\textbf{In the family, and next door}

\noindent
\begin{tabularx}{\linewidth}{@{}>{\raggedright\arraybackslash}X>{\raggedright\arraybackslash}X>{\raggedright\arraybackslash}X@{}}
\toprule
\textbf{} & \textbf{word} & \textbf{same root} \\
\midrule
Kannada & \ka{ಸಂತೆ} (\emph{sante}) & yes \\
Telugu & \te{సంత} (\emph{santa}) & yes \\
Malayalam & \ml{ചന്ത} (\emph{canta}) & yes \\
Hindi (neighbour) & \dv{बाज़ार} (\emph{bāzār}) &  \\
English & a market &  \\
\bottomrule
\end{tabularx}

\begin{grammarlens}[title={What you've built}]
One more word for this chapter.
\end{grammarlens}

\subsection*{Your turn}
\begin{itemize}
\raggedright
  \item \emph{Say it:} \emph{cantai}
  \item \emph{Say it:} \emph{cantai}, once more
  \item \emph{Say it:} say \emph{cantai}
  \item \emph{Recall:} say the Tamil for a library, then the Tamil for a bank, then say \emph{cantai} again
\end{itemize}

\begin{tcolorbox}[breakable,colback=teal!4,colframe=teal!35!black,title={Before you move on}]
What does \ta{சந்தை} mean? (\textquotedblleft{}A market\textquotedblright{}.) Say it once more.
\end{tcolorbox}
\section[\texorpdfstring{\ta{கடற்கரை} (kaṭaṟkarai)}{கடற்கரை (kaṭaṟkarai)}]{\ta{கடற்கரை} (kaṭaṟkarai) — a beach}

\label{lesson:TA-C115-katarkarai}



\begin{quote}
Before the new one: say the Tamil for a bank, then the Tamil for a market.
\end{quote}

\subsection*{You'll want to know: \ta{கடற்கரை}}
\textbf{\ta{கடற்கரை}} — \emph{kaṭaṟkarai} — \textquotedblleft{}a beach\textquotedblright{}.

\textbf{In the family, and next door}

\noindent
\begin{tabularx}{\linewidth}{@{}>{\raggedright\arraybackslash}X>{\raggedright\arraybackslash}X@{}}
\toprule
\textbf{} & \textbf{word} \\
\midrule
Kannada & \ka{ಕಡಲತೀರ} (\emph{kaḍalatīra}) \\
Malayalam & \ml{കടപ്പുറം} (\emph{kaṭappuṟaṁ}) \\
English & a beach \\
\bottomrule
\end{tabularx}

\begin{grammarlens}[title={What you've built}]
One more word for this chapter.
\end{grammarlens}

\subsection*{Your turn}
\begin{itemize}
\raggedright
  \item \emph{Say it:} \emph{kaṭaṟkarai}
  \item \emph{Say it:} \emph{kaṭaṟkarai}, once more
  \item \emph{Say it:} say \emph{kaṭaṟkarai}
  \item \emph{Recall:} say the Tamil for a bank, then the Tamil for a market, then say \emph{kaṭaṟkarai} again
\end{itemize}

\begin{tcolorbox}[breakable,colback=teal!4,colframe=teal!35!black,title={Before you move on}]
What does \ta{கடற்கரை} mean? (\textquotedblleft{}A beach\textquotedblright{}.) Say it once more.
\end{tcolorbox}
\section[\texorpdfstring{\ta{அலுவலகம்} (aluvalakam)}{அலுவலகம் (aluvalakam)}]{\ta{அலுவலகம்} (aluvalakam) — an office}

\label{lesson:TA-C115-aluvalakam}



\begin{quote}
Before the new one: say the Tamil for a market, then the Tamil for a beach.
\end{quote}

\subsection*{You'll want to know: \ta{அலுவலகம்}}
\textbf{\ta{அலுவலகம்}} — \emph{aluvalakam} — \textquotedblleft{}an office\textquotedblright{}.

\textbf{In the family, and next door}

\noindent
\begin{tabularx}{\linewidth}{@{}>{\raggedright\arraybackslash}X>{\raggedright\arraybackslash}X@{}}
\toprule
\textbf{} & \textbf{word} \\
\midrule
Kannada & \ka{ಕಚೇರಿ} (\emph{kacēri}) \\
Telugu & \te{కార్యాలయం} (\emph{kāryālayaṁ}) \\
Malayalam & \ml{ഓഫീസ്} (\emph{ōphīs}) \\
Hindi (neighbour) & \dv{दफ़्तर} (\emph{daftar}) \\
English & an office \\
\bottomrule
\end{tabularx}

\begin{grammarlens}[title={What you've built}]
One more word for this chapter.
\end{grammarlens}

\subsection*{Your turn}
\begin{itemize}
\raggedright
  \item \emph{Say it:} \emph{aluvalakam}
  \item \emph{Say it:} \emph{aluvalakam}, once more
  \item \emph{Say it:} say \emph{aluvalakam}
  \item \emph{Recall:} say the Tamil for a market, then the Tamil for a beach, then say \emph{aluvalakam} again
\end{itemize}

\begin{tcolorbox}[breakable,colback=teal!4,colframe=teal!35!black,title={Before you move on}]
What does \ta{அலுவலகம்} mean? (\textquotedblleft{}An office\textquotedblright{}.) Say it once more.
\end{tcolorbox}
